\section{Profile spaces and corrected test states}
\label{lift:section}

Throughout this section the mass $m\geq0$ is fixed. Projection to the
fast line identifies where a bounded family can concentrate. To identify
its limiting charge, we also need test states
whose charge outputs converge in norm. The uncorrected oscillator
product has an order-one component perpendicular to the fast line.
The inverse of the fast charge removes this component with a correction
of size $h^{-1/2}$. Such reduced-inverse corrections occur in the
asymptotic expansions of Konechny~\cite{Konechny1999},
Fr\"ohlich et al.~\cite{FroehlichEtAl2000}, and
Lin--Yin~\cite[Section~4.1]{LinYin2015}. Here we prove convergence of
the corrected charge output in the norm required for duality. We give this construction without imposing any
moment condition on an internal $L^2$ kernel vector.

\begin{definition}[Profile space]\label{def:profile-space}
For an unordered partition $C=(n_1,\ldots,n_k)$ of $N$, with $k\geq2$,
let its labeled center space be
\[
 \mathcal Z_C=\left\{(c_1,\ldots,c_k)\in(\R^9)^k:
                         \sum_\alpha n_\alpha c_\alpha=0\right\},
 \qquad
 |\dd c|^2=\sum_\alpha n_\alpha|\dd c_\alpha|^2.
\]
Write $\mathcal Z_C^\circ$ for the complement of the coincidence
sets $c_\alpha=c_\beta$. Its center Clifford module is denoted by
$\mathcal F_C$. Set $\mathcal K_n=\Ker H_n$ and
$\mathcal K_1=\C$, and form the Hilbert tensor product
\[
 \mathcal V_C=\mathcal F_C\,\widehat\otimes\,
                        \widehat\bigotimes_{\alpha=1}^k\mathcal K_{n_\alpha}.
\]
The profile Hilbert space $\mathcal P_C$ is
$L^2(\mathcal Z_C^\circ;\mathcal V_C)$ with the simultaneous
coordinate and fiber symmetry for permutations of equal-sized blocks.
The fiber action uses the natural center Clifford action and the
graded permutation of internal factors. Lebesgue measure is that of
the displayed trace metric, divided by the number of such
permutations; equivalently one uses its quotient measure. No
finite-dimensionality of the spaces $\mathcal K_n$ is built into
this definition.

On $\mathcal P_C$ the candidate charge is the abelian center operator
\begin{equation}
 D_v^C(m)=\theta_C(A_i)\cdot p_{c_i}
                          +m\theta_C(L_i)\cdot c_i,
 \label{lift:center-charge}
\end{equation}
where all contractions are in orthonormal trace coordinates on
$\mathcal Z_C$. Define $\mathcal T_C$ to be the space of finite sums
of smooth compactly supported center functions on
$\mathcal Z_C^\circ$ with fiber values in algebraic tensor products
of actual internal kernel vectors, symmetrized as above.
\end{definition}

The tests $\mathcal T_C$ are dense in $\mathcal P_C$ even before
anything is known about the dimensions of its internal factors:
finite algebraic tensors are dense in a Hilbert tensor product,
ordinary compactly supported smooth functions are dense in $L^2$ on
the open center space, and the finite permutation average is a
bounded projection. The graph closure across collisions will be
established in Lemma~\ref{lem:collision-removal}; at this stage the
tests are deliberately supported away from collisions.

We use the following notation for their leading identification. In
a chart of Section~\ref{fast:section}, let $c=c(b)$ and $\xi=b-c$.
The internal variables have total dimension
\[
 d_{\mathrm{int}}=9\sum_{\alpha=1}^k(n_\alpha^2-1).
\]
For $f\in\mathcal T_C$, interpret each kernel factor as its wavefunction
in $y=h^{1/3}\xi$, and multiply the resulting unitary rescaling by
fixed smooth block-invariant cutoffs in $\xi$, equal to one near zero
and supported inside the chart. Denote the resulting slow state by
$f_h$. The unitary internal scaling contributes the factor
$h^{d_{\mathrm{int}}/6}$. The leading local identification is
\begin{equation}
 \mathcal J_hf=\chi_v(c)f_h
 \label{lift:leading-identification}
\end{equation}
in the flat $(b,t)$ coordinates; the center section $\chi_v(c)$ is
independent of $v$. When necessary a fixed physical $z$ cutoff,
equal to one near zero, is included. Different fixed cutoffs of
this kind change \eqref{lift:leading-identification} by $o(1)$ in
norm. We will show below that chart choices are equally immaterial.

\begin{proposition}[Corrected test lifts]\label{prop:test-lifts}
Fix $m\geq0$, a unit spinor $v$, and a partition $C$ with $k\geq2$.
For every $f\in\mathcal T_C$ there are gauge-singlet graph-domain
states $I_{h,v,m}f$, defined for all sufficiently large $h$, such that
\begin{align}
 \|I_{h,v,m}f-\mathcal J_hf\|&\longrightarrow0,\notag\\
 \|D_v(h,m)I_{h,v,m}f-
                  \mathcal J_h(D_v^C(m)f)\|&\longrightarrow0.
 \label{lift:graph-convergence}
\end{align}
The construction is linear on any fixed finite-dimensional space
of tests, and both conclusions hold simultaneously for every fixed
finite test family. In particular, all inner products of their
leading states and indicated charge outputs converge to the
corresponding profile inner products. For distinct partitions the
families, including their outputs, are asymptotically orthogonal.
The leading identifications intertwine block permutations and
$\Spin(3)\times\Spin(6)$; each fixed symmetry may be passed through
these norm limits. The lifts may depend on $v$ and $m$.
\end{proposition}

\begin{proof}
We first work in one labeled center chart. All its cutoffs may be
chosen the same for a fixed finite family of tests. Internal kernel
vectors have the local Sobolev regularity supplied by
Lemma~\ref{lem:domains}. Thus the expressions below are distributionally
defined with the indicated local derivatives, even when a kernel
vector has no global moment or global derivative bound.

Consider one internal factor $\phi\in\mathcal K_n$. If its physical
cutoff is $\kappa(\xi)$, its unitary rescaling is
\[
 \phi_h(\xi)=h^{9(n^2-1)/6}\phi(h^{1/3}\xi).
\]
Exact scaling of the undeformed charge gives
$D_v(h,0)\phi_h=0$ distributionally. Hence
\[
 D_v(h,0)(\kappa\phi_h)
       =-\ii\theta(A_i)\cdot(\partial_i\kappa)\phi_h.
\]
The derivative of $\kappa$ is bounded and supported at a fixed
positive physical distance from zero. Changing variables to $y$
shows that its right-hand side tends to zero in norm by the $L^2$
tail of $\phi$. This holds for every internal factor and every
fixed spinor. It also shows that the cutoff products converge in
norm, under internal rescaling, to the uncut products.

The crucial additional estimate is separate smallness of the
derivative and bracket terms:
\begin{equation}
 \|p_b f_h/\sqrt h\|
       +\sum_{i<j}\|\sqrt h\,[b_i,b_j]f_h\|
          \longrightarrow0.
 \label{eq:kernel-scaled-estimate}
\end{equation}
To prove it, use the undeformed square identity on each cut-off
internal factor and divide by $h$. Its fermion term is bounded in
absolute value by $Ch|\xi|$. Summing the internal factors gives
\begin{equation}
 \frac1h\|p_\xi f_h\|^2
 +h\sum_{i<j}\|[b_i,b_j]f_h\|^2
 \leq\frac C h\sum_\alpha
           \|D_{v,\alpha}(h,0)f_h\|^2
           +C\int\!\Bigl(\sum_\alpha|\xi_\alpha|\Bigr)|f_h|^2.
 \label{lift:separate-energy}
\end{equation}
Here the bracket is internal, since the scalar centers commute
with all block variables. The first term tends to zero by the
cutoff calculation. In the second, every $\xi_\alpha$ is bounded
on the fixed physical support, while its norm measure concentrates
at zero. To verify convergence directly, split its integral at
$|\xi_\alpha|=\varepsilon$: the inner contribution is bounded by
$C\varepsilon$, and the outer contribution tends to zero by the
rescaled $L^2$ tail. First let $h\to\infty$, then
$\varepsilon\downarrow0$. Center derivatives of $f_h$ are uniformly
bounded, so adding them divided by $\sqrt h$ proves
\eqref{eq:kernel-scaled-estimate}. The same splitting proves
$\|\xi f_h\|\to0$; the fixed-mass internal term is therefore
negligible. This proof uses only bounded physical cutoffs and
$L^2$ of the uncut kernels.

We now identify the output that the correction must cancel. Let
$R_v^{(0)}(b)$ be the density-conjugated remainder in
\eqref{fast:R-coordinate}--\eqref{fast:density-term} after removing
the fast mass term and the term containing $h^{-1/2}p_b$.
Thus $R_v^{(0)}$ contains the quadratic $[t,t]$ multiplier, the
$tp_t$ term, the orbit-spin term, and the density term. It is
independent of $h,m$ and differentiates only the fast variables.
Define the oscillator-vector map on the slow Clifford fiber by
\begin{equation}
 A(b)=-\ii\sum_{i,Q}\theta_s^Q(A_i)
                         \partial_{b_i^Q}\chi_v(b)
                       +R_v^{(0)}(b)\chi_v(b).
 \label{lift:residual-map}
\end{equation}
The fast output of the uncorrected state $\chi_v(b)f_h$ is zero,
and its full charge output is
\begin{equation}
 D_v(h,m)(\chi_v(b)f_h)
 =\chi_v(b)(D_v^C(m)f)_h+A(b)f_h+o_{L^2}(1).
 \label{lift:uncorrected-output}
\end{equation}
Indeed the internal undeformed outputs and internal mass terms
tend to zero by the cutoff calculation above. The removed slow
derivative term tends to zero by
\eqref{eq:kernel-scaled-estimate} when it differentiates $f_h$,
and by the bounded derivatives of $\chi_v$ when it differentiates
the section. The fast mass term is $O(m/\sqrt h)$ on these states.
Lemma~\ref{lem:fast-line} also gives boundedness of $A(b)$ and its
required slow derivatives in every prescribed finite oscillator
Sobolev order.

Lemma~\ref{lem:diagonal-compression} says $P_v(c)A(c)=0$. Smoothness
and concentration of $f_h$ give
\[
 \|P_v(b)A(b)f_h\|\longrightarrow0.
\]
Define the remaining perpendicular map and its correction by
\begin{equation}
 \begin{split}
 V(b)&=(1-P_v(b))A(b),\qquad
 W(b)=\bigl(E_v(b)|_{\Ran(1-P_v(b))}\bigr)^{-1}V(b),\\
 I_{h,v,m}f&=\chi_v(b)f_h-h^{-1/2}W(b)f_h.
 \end{split}
 \label{lift:correction}
\end{equation}
The inverse here acts with the Clifford sign it has on the full
graded module. Thus $E_v(b)W(b)=V(b)$ exactly. The maps preserve
fast $Z_0$ neutrality, since every term in the construction is
equivariant. Also $W(b)$ is even as a map of the full graded
fibers: $V$ is odd and the inverse of $E_v$ is odd. The correction
therefore preserves the parity of a test.

We verify that all other outputs of the correction tend to zero.
For any smooth oscillator-vector map $M(b)$ whose coefficients
and first slow derivatives are bounded in the requisite oscillator
norms, the explicit slow charge gives
\begin{align}
 \|h^{-1/2}D_v^s(h,m)(M f_h)\|
 \leq C_M\bigl(&\|p_b f_h/\sqrt h\|
      +\sum_{i<j}\|\sqrt h[b_i,b_j]f_h\|\notag\\
      &+h^{-1/2}(1+m)\|f_h\|\bigr).
 \label{lift:matrix-stability}
\end{align}
The boundedness of $b$ on the support is included in $C_M$.
The estimate follows by applying each derivative to $M$ or $f_h$
and by bounding the remaining finite Clifford matrices. In
particular it remains true when such matrices act on internal
kernel factors. There is no assertion that those matrices
preserve the internal kernels. Equation~\eqref{eq:kernel-scaled-estimate}
is precisely what makes the right-hand side tend to zero.

Taking $M=W$ proves smallness of the slow output of the correction.
For its remainder output, \eqref{fast:R-coordinate} gives
\[
 \|h^{-1/2}R_v(W f_h)\|
 \leq C\bigl(h^{-1/2}(1+m/\sqrt h)\|f_h\|
                    +h^{-1}\|p_b f_h\|\bigr)=o(1).
\]
All $t$ multiplications and $t$ derivatives here act on $W$ and
are bounded by its oscillator norms. Finally, the fast output
of the correction is $-V(b)f_h$, exactly canceling the remaining
uncorrected output. We have proved
\[
 D_v(h,m)I_{h,v,m}f
                  =\chi_v(b)(D_v^C(m)f)_h+o(1).
\]
The correction has norm $O(h^{-1/2})$.
Replacing $\chi_v(b)$ by $\chi_v(c)$ costs $o(1)$, both for $f$
and for the fixed test $D_v^C(m)f$, by continuity and concentration.
This proves \eqref{lift:graph-convergence} in the chart.

The states can be taken in the actual graph domain. Their slow
supports lie in the interior of the chosen spectral-window chart,
and their fast factors have bounded oscillator moments of every
fixed order used above. If a physical transverse cutoff
$\tau(z)=\tau(t/\sqrt h)$ is desired, its deleted tail and charge
commutator tend to zero: on that tail $|t|\geq a\sqrt h$, whereas
arbitrarily high fixed oscillator moments are bounded. The
ambient charge commutator with a fixed smooth physical cutoff is
a bounded Clifford multiplier. Extending the resulting function
by gauge equivariance and by zero outside the chart introduces
no boundary jump. The computed distributional charge is in $L^2$,
so Lemma~\ref{lem:domains} puts the state in the graph domain.
Local smoothing, followed by compact-group averaging if required,
gives core approximations with the same graph limits.

The remaining global and symmetry assertions follow from the
chart compatibility proved next. In particular a finite partition
of unity on the compact center support produces the asserted
global lift and preserves linearity on a fixed finite test space.
\end{proof}

\begin{lemma}[Compatibility of leading identifications]
\label{lem:chart-compatibility}
On overlaps of separated-center charts, the identifications
$\mathcal J_h$ agree to $o(1)$ in norm for each fixed test, with
the natural center and internal fiber transports. The same is
true for the indicated outputs $\mathcal J_hD_v^C(m)f$. A change
of separating axis shifts the block variable $b$ by $O(|z|^2)$;
its leading fast coordinate change is projection along the
central conjugation directions, with the positive unitary
density factor. The statements hold for finite test families,
block relabelings, and each fixed spatial symmetry in
$\Spin(3)\times\Spin(6)$.
\end{lemma}

\begin{proof}
Choose two separating axes valid on the overlapping compact
center support. To put $b+z$ into the second slice, solve its
off-block slice equation by a gauge rotation with off-block
generator $w(b,z)$. The differential of that equation at $z=0$
is the invertible map $w\mapsto[w,b_{e'}]$ on $\mathfrak n$.
The implicit function theorem with the fixed spectral gap gives
$w(b,z)=O(|z|)$ with uniformly bounded smooth coefficients.
Since $[\mathfrak n,\mathfrak k_0]\subset\mathfrak n$, its first
action on $b$ is off-block. Thus the block-diagonal part of the
transformed configuration satisfies
\begin{equation}
 b'=b+O(|z|^2),
 \qquad
 z'=L_bz+O(|z|^2),
 \label{lift:chart-transition}
\end{equation}
where $L_c$ is the linear change of slice obtained by projecting
along $([w,c_i])_i$. A fixed change of block frames or labels can
be applied before this calculation.

On bounded rescaled sets $|y|,|t|\leq M$,
\[
 h^{1/3}|b'-b|=O_M(h^{-2/3}),\qquad
 t'=L_ct+o(1),\qquad c'=c+O_M(h^{-1}).
\]
The spin action of the gauge rotation tends uniformly to the
identity there. The density factors in
\eqref{fast:density}, together with the Jacobian of $L_c$, are
exactly those of the unitary change between the two flat-measure
transverse slices. This can also be seen by mapping each slice
to the same perpendicular slice for $([w,c_i])_i$: both start
from the same ambient Euclidean volume and the same homogeneous
space orbit measure. All these density square roots are
positive. The transformed Gaussian is therefore the same
normalized oscillator Gaussian, and its spin factor agrees by
the phase convention of Lemma~\ref{lem:fast-line}.

For bounded smooth internal profiles these statements prove norm
convergence by change of variables and dominated convergence on
bounded $(y,t)$ sets. They extend to an arbitrary fixed internal
$L^2$ kernel tensor by approximation: truncate and approximate
that tensor on a bounded $y$ set, use the unitary coordinate
changes and bounded physical cutoffs to bound the approximation
error, and then remove the truncation. The fast tails are
uniformly small by the oscillator moments. No estimate of an
uncut $y$ moment is involved. The same argument applies to
$D_v^C(m)f$, which has the same finite collection of internal
kernel factors and a smooth compact center coefficient.

For block relabelings the only additional transformation is the
ordinary graded transport of the slow factors: the fast line
has no extra character by Lemma~\ref{lem:fast-line}. Rotating a
chart by a fixed element $g$ of $\Spin(3)\times\Spin(6)$ gives
the same construction and the natural profile action. Its
charge output transforms with $v$ replaced by $gv$, as follows
from charge covariance. In particular for rotations fixing
$v$ it transforms the same charge. These are statements for
each fixed $g$; a uniform assertion over all group elements is
not needed.

To finish the global construction, take smooth functions
$\lambda_\nu(c)$ subordinate to finitely many labeled center
charts and summing to one on the test support. Sum the local
lifts of $\lambda_\nu f$, making the finite block permutation
average when passing to the quotient. The overlap calculation
identifies their leading sum with $\mathcal J_hf$, and their
outputs with
\[
 \mathcal J_h\sum_\nu D_v^C(m)(\lambda_\nu f)
                         =\mathcal J_hD_v^C(m)f.
\]
The product-rule terms cancel because
$\sum_\nu\dd\lambda_\nu=0$ on the support. Thus the partition
introduces exactly the center Dirac product rule already
present in the required output.

Finally, compact center supports away from collisions for
distinct unordered partitions are disjoint compact subsets
of the commuting gauge quotient. The norm measures of their
leading products concentrate on those supports. Disjoint
neighborhoods and Cauchy--Schwarz therefore make their mutual
inner products tend to zero. Charge outputs have the same
limiting supports, so this also proves orthogonality of the
indicated outputs and the mixed inner products. Within one
partition, the unitary rescalings, normalized fast vector, and
fixed orbit-volume convention give exactly the profile inner
products; polarization proves the finite-family assertion.
\end{proof}

For later symmetry averaging, one useful consequence deserves
to be explicit. If a profile test is invariant under a compact
group of the stated symmetries, its lift and each fixed rotated
lift have difference tending to zero in norm. Haar averaging
has the same leading profile by dominated convergence, because
their norms are bounded. If the group commutes with the selected
charge, averaging does not increase the charge error. This uses
only the pointwise compatibility in Lemma~\ref{lem:chart-compatibility}.

The compatible tests identify limiting center equations. We next prove
that their profiles capture all norm, including norm at internal scales
that shrink in physical coordinates.
